% !TeX root=../main.tex
% در این فایل، عنوان پایان‌نامه، مشخصات خود و چکیده پایان‌نامه را به انگلیسی، وارد کنید.

%%%%%%%%%%%%%%%%%%%%%%%%%%%%%%%%%%%%
\latinuniversity{K. N. Toosi University of Technology}
\latincollege{Faculty of Computer Engineering}
\latinfaculty{Faculty of Computer Engineering}
\latindepartment{Artificial Intelligence Group}
\latinsubject{Computer Engineering}
%\latinfield{field}
\latintitle{Unsupervised Depth Estimation Using Internal Calibration Parameters as Input}
\firstlatinsupervisor{Dr. Behrouz Nasihatkon}
%\secondlatinsupervisor{}
%\firstlatinadvisor{}
%\secondlatinadvisor{}
\latinname{Mohammad Amin}
\latinsurname{Ghasemzadeh}
\latinthesisdate{Autumn 2026}
\latinkeywords{
Monocular Depth Estimation,
Unsupervised Deep Learning,
Metric Depth,
Camera Intrinsics,
Imaging Geometry,
Depth Anything V2
}

\enabstract{
Monocular depth estimation is intrinsically ambiguous because a single image does not uniquely determine the metric scale of a scene. This thesis investigates whether supplying camera intrinsics explicitly can reduce that ambiguity in a foundation depth model. An extension of Depth Anything V2 is developed in which a normalized nine-dimensional camera descriptor is encoded as a gated token and injected into the DINOv2 encoder. The experimental methodology was designed and evaluated through a sequence of controlled investigations that addressed geometrically inconsistent crop augmentation and a silent metric-head mismatch during evaluation.

All final comparisons use the standard 654-image NYUv2 Eigen test set and two matched training arms: the proposed model with the camera calibration matrix and the control model without the camera calibration matrix. The remainder of the thesis uses the names proposed model and control model without repeating these input conditions. Under relative evaluation, the proposed ViT-S model improves over both the released model and the matched control, while the proposed ViT-L model does not show a significant contribution because per-image scale-and-shift alignment removes the principal information supplied by the focal length. Under metric evaluation without alignment, the proposed ViT-L model consistently improves over the matched control, and the effect grows as field-of-view diversity increases. A geometrically exact focal-jitter augmentation is then introduced to make scale unrecoverable from image pixels alone. With this augmentation, the proposed model obtains AbsRel 0.0849 versus 0.1030 for the control and improves every reported metric with Wilcoxon significance of at least $p\leq10^{-24}$.

Mechanistic and causal tests confirm that the improvement is genuinely attributable to the camera calibration matrix. The log--log slope of predicted depth versus injected focal scale rises from 0.06 before focal jitter to 0.93 to 0.97 afterward, close to the pinhole-model ideal of one. Withholding the camera calibration matrix increases AbsRel by 59 percent, while deliberately corrupting focal length produces errors of the magnitude predicted by projective geometry. These results show that camera intrinsics can provide a statistically significant and causally necessary scale signal, while also revealing sensitivity to calibration accuracy and the need for future evaluation on genuinely multi-camera and outdoor data.
}
